\begin{claim}{D.45}{Proved}{A characteristic binary-icosahedral refinement}
Let $F=\widehat F(x,y)$ and let $R$ be D.38's characteristic
kernel with $F/R=A_5^{19}$.  Intersecting the kernels of all
epimorphisms $F\twoheadrightarrow S=\mathrm{SL}_2(\F_5)$ gives
an open characteristic subgroup $\widetilde R\subset R$ and
the following exact quotient and reduction map:
\[
 \widetilde Q=F/\widetilde R
       \xrightarrow{\ \sim\ }S^{19}\times C_2^2,
 \qquad
 \pi:\widetilde Q\longrightarrow A_5^{19},\quad
       ((g_i),a)\longmapsto(g_i\bmod\{\pm I\})_i.
                                                        \tag{D.45a}
\]
For the $10\,125\,000$ finite tuples $t$ of D.38, impose at
this specified refinement the normalized multiplier $-1$
action, both Harbater--Schneps cycle equations, and the
exact norm $f\alpha(f)=1$.
Exactly $5\,062\,500$ base tuples admit such a refined
$f$-coordinate.  They are precisely D.38's cusp-triangle
survivors, and each has $64$ admissible $f$-coordinates.
Thus the total number of refined coordinate pairs is
$324\,000\,000$.
For a fixed $t$, these $64$ coordinates define the same
actual automorphism of $\widetilde Q$, of order two.
In particular this named characteristic refinement
detects exactly D.38's first excluded half.  All
$2\,531\,250$ triangle survivors outside D.38's finite
pure conjugacy set still pass this refinement.
\end{claim}
\begin{proof}
All matrix calculations below are in $\F_5$.
The program enumerates the $120$ determinant-one matrices
and verifies that their commutator subgroup has order $120$.
Their center is $\{I,-I\}$, and their projective quotient is
identified explicitly with D.38's $A_5$ by a generating pair
of orders two and three with product of order five.
The resulting bijection is checked against all $60^2$
products.  Conjugation by the $120$ projective invertible
matrices gives all automorphisms of $S$.
Indeed $\Aut(S)\to\Aut(A_5)$ is injective: an element of
its kernel has the form $g\mapsto c(g)g$, where
$c:S\to\{\pm I\}$ is a homomorphism, and perfection makes
$c=1$.  D.38 gives $|\Aut(A_5)|=120$, and the enumerated
projective conjugations already give $120$ distinct
automorphisms of $S$.

Choose lifts $X_i,Y_i\in S$ of the nineteen generating
pairs in D.38, and put $Z_i=(X_iY_i)^{-1}$.
Every pair of lifts generates $S$: its image generates
$A_5$, so a proper image would have index two in $S$,
contrary to perfection.  The program also checks all
$4\cdot19=76$ choices directly.
Two maps to $A_5$ with a given kernel differ by an
automorphism, and every automorphism lifts uniquely to
$S$.  Over each of the nineteen $A_5$ kernels there
are exactly four $S$ kernels, from the independent
sign choices for the lifts of $x$ and $y$.
If two such lifts had the same kernel, the induced
automorphism of $S$ would act trivially on $A_5$;
the proved injectivity would identify the lifts.
Consequently these are exactly the $76$ epimorphism
kernels defining $\widetilde R$.

The joint map
\[
 x\longmapsto((X_i),e_x),\qquad
 y\longmapsto((Y_i),e_y),
        \qquad e_x,e_y\text{ the basis of }C_2^2,
                                                        \tag{D.45b}
\]
is onto.  To see this without a rank inference, its
projection $D$ to $S^{19}$ maps onto $A_5^{19}$ by
D.38.  Thus $D Z(S^{19})=S^{19}$; centrality gives
$[S^{19},S^{19}]=[D,D]\leq D$, so perfection gives
$D=S^{19}$.  The commutator subgroup of the image
of (D.45b) therefore contains $S^{19}\times\{0\}$,
and its projection to $C_2^2$ is onto.  This proves
surjectivity of (D.45b).
Every other sign lift of the $i$th pair is obtained by
composing (D.45b) with
\[
 ((g_j),a)\longmapsto g_i(-I)^{\ell(a)},
                         \qquad \ell\in\Hom(C_2^2,C_2).
\]
The intersection of these four kernels, for all $i$,
is trivial: the zero character first makes each
$g_i=I$, and all characters then make $a=0$.
This proves (D.45a), including the specified kernel.
It is characteristic because every automorphism of
$F$ permutes its epimorphisms to $S$.

Here are the complete finite equations used in the
refinement.  The named word substitutions $\theta,
\omega$ of D.38 descend to $\widetilde Q$.
On its characteristic perfect subgroup $S^{19}$ they
permute the factors by D.38's recorded maps and apply
the unique lifted $S$ automorphisms just described.
Write $\widetilde t_i\in S$ for any lift of $t_i$ and
put
\[
 f_{0,i}=\theta(\widetilde t)_i\widetilde t_i^{-1},
 \qquad
 f_i=(-I)^{b_i}f_{0,i},\qquad b_i\in\F_2.
                                                        \tag{D.45c}
\]
These are all the $2^{19}$ lifts of
$f_Q=\theta(t)t^{-1}$ into
$\widetilde Q'=S^{19}$.
For every such lift, set
\[
 \begin{split}
 c_{x,i}&=X_i^{-1}
       (\widetilde t_i X_i\widetilde t_i^{-1})^{-1},\\
 c_{y,i}&=(f_i^{-1}Y_i^{-1}f_i)
       (\widetilde t_i Y_i\widetilde t_i^{-1})^{-1}.
 \end{split}                                             \tag{D.45d}
\]
D.38's normalization implies $c_{x,i},c_{y,i}\in
\{I,-I\}$.  The actual automorphism induced by the
normalized pair is exactly
\[
 \alpha((g_i),(a_x,a_y))=
 \bigl((\widetilde t_i g_i\widetilde t_i^{-1}
           c_{x,i}^{a_x}c_{y,i}^{a_y})_i,
                      (a_x,a_y)\bigr).
                                                        \tag{D.45e}
\]
This formula is a continuous automorphism; its
restriction to the first factor is an automorphism,
and its two correction characters are central.
It sends the two generators (D.45b) to the required
$X^{-1}$ and $f^{-1}Y^{-1}f$.
Since $\widetilde t_i^2$ is central and both correction
characters have exponent two, its square is the
identity.  Its reduction is the nonidentity order-two
automorphism of D.38, so its order is exactly two.
Multiplying $f_i$ by a central sign changes neither
(D.45d) nor (D.45e).  Conversely these images of the
two generators determine the automorphism uniquely.

The program tests the following three exact equations
in $S^{19}$ for all choices (D.45c):
\[
 \theta(f)f=1,\qquad
 \omega^2(f)Z^{-1}\omega(f)Y^{-1}fX^{-1}=1,
 \qquad f\alpha(f)=1.
                                                        \tag{D.45f}
\]
On $f$, which has zero $C_2^2$ coordinate, formula
(D.45e) is conjugation by $\widetilde t$.
Thus every operation in (D.45f) is a specified
matrix multiplication or a specified factor map;
the norm is tested as an equation, not inferred
from the order of the finite automorphism.

The factor maps for $\theta,\omega$ have six joint
orbits.  Exhaustion within these orbits gives the
following complete counts, with D.38's zero-based
factor indices:
\[
\begin{array}{c|r|l}
 \text{factor orbit}&\text{base tuples}&
       \text{numbers of admissible central lifts}\\ \hline
 \{0,2,10\}&15&2\text{ for each of }15\\
 \{1,9,17\}&15&2\text{ for each of }15\\
 \{3,5,6,7,11,14\}&60&0\text{ for }30;\quad4\text{ for }30\\
 \{4,12,18\}&15&2\text{ for each of }15\\
 \{8,13,15\}&10&2\text{ for each of }10\\
 \{16\}&5&1\text{ for each of }5.
\end{array}                                               \tag{D.45g}
\]
The rejected thirty assignments in the six-factor
orbit are exactly its triangle failures; the other
five orbits have no triangle failure.
The equations on one joint orbit involve no other
orbit, so their solution sets have the displayed
Cartesian-product decomposition.
The surviving count is
$15\cdot15\cdot30\cdot15\cdot10\cdot5=5\,062\,500$,
and the number of central lifts per survivor is
$2\cdot2\cdot4\cdot2\cdot2\cdot1=64$.
This proves every asserted count.
The executed implementation is
\path{code/part_a/a5_schur_refinement.py}; its output
records each orbit, the lift histogram, and sample
coordinates.  These computations impose a necessary
finite refinement on full involutions.  They assert
no full lift of a surviving finite coordinate pair.
\end{proof}
